\chapter{Training and evaluation}

\section{From raw data to results}

Figure~\ref{fig:training} follows the data through the pipeline. The LLVIP~\cite{jia2021llvip} and
M3FD~\cite{liu2022tardal} archives are converted once into one paired layout: two folders of images with identical
names, the long side at 640 pixels, and labels in YOLO format. A paired loader reads both images and stacks them
into one four-channel tensor, visible colour plus thermal grey. Augmentation is shared, so a crop or a flip moves
both cameras together, except for colour jitter, which touches only the visible image. Modality dropout then blacks
out one camera with probability 0.15 each, and in half of those cases tells the model which camera is gone, so that
both the silent and the reported failure are seen in training.

\widefig{fig_training}{The training and evaluation pipeline. Brown: data; purple: training; gold: evaluation;
blue: what the project delivers. The loop between the losses and the optimiser is one training step.}{fig:training}

Training follows YOLO26: two task-aligned losses, one-to-many with ten predictions per object and one-to-one with a
single one, summed. The optimiser is chosen by Ultralytics and was AdamW at this schedule. Mixed precision runs the
network in 16-bit floats, but the whole scan path, from the projections to the fusion weights, stays in 32-bit
floats, because a recurrence over twelve thousand tokens overflows the 16-bit range. Two T4 GPUs share each batch
through distributed data parallelism.

Every trained checkpoint goes through the same evaluation. COCO average precision is computed on the full test
set, with the AI-TOD size bands that separate very tiny, tiny and small objects. A degradation probe applies seven
sensor faults at test time, from a darkened visible image to a thermal camera shifted by sixteen pixels. Latency is
measured on one T4. The five hypotheses of the pilot are then decided by rules written down before training.

\section{The notebooks}

The work is organised as notebooks, each one step of Figure~\ref{fig:workflow}. Phase~0 sets up and checks the
environment and prepares the data. Phase~1 trains the single-camera and concat baselines, \method{}, its ablations
and the head control in the pilot, then the improved version in two further rounds, and closes with the probe, the
results and a showcase of the trained models.

\widefig{fig_workflow}{The notebooks and how their outputs flow. Each training notebook runs as an independent
Kaggle session; the probe reads their checkpoints, and the results notebook decides the hypotheses.}{fig:workflow}
